\documentclass{article}
\usepackage[T1]{fontenc}
\usepackage{lmodern}
\usepackage{graphicx}
\usepackage{amsmath,amssymb}
\usepackage[a4paper,margin=2cm]{geometry}
\usepackage[absolute,overlay]{textpos}
\setlength{\TPHorizModule}{1mm}
\setlength{\TPVertModule}{1mm}
\usepackage[indonesian]{babel}
\usepackage{hyperref}
\setlength{\emergencystretch}{2em}
% unit: Halaman 54

\begin{document}

% === Halaman 54 (llm) ===

\begin{itemize}
  \item Contoh:
  \vspace{5mm}
  \textcolor{red}{DYDUXRMH}TVDV\textcolor{cyan}{NQD}QNW\textcolor{red}{DYDUXRMH}ARTJGW\textcolor{cyan}{NQD}
\end{itemize}

Kriptogram yang berulang: \textbf{DYUDUXRMH} dan \textbf{NQD}.

Jarak antara dua buah perulangan \textbf{DYUDUXRMH} $= 18$.
Semua faktor pembagi $18 : \{18, 9, 6, 3, 2\}$

Jarak antara dua buah perulangan \textbf{NQD} $= 20$.
Semua faktor pembagi $20 : \{20, 10, 5, 4, 2\}$.

Irisan dari kedua buah himpunan tersebut adalah 2
Panjang kunci kemungkinan besar adalah 2.

\end{document}
